\centerline{\bf \Huge Сложный тургор}

\begin{flushright}
        \scriptsize{Все задачи легко гуглятся, поэтому исключительно на вашу честность!}
\end{flushright}

\problem{1}
Десяти ребятам положили в тарелки по 100 макаронин. Есть ребята не хотели и стали играть. Одним действием кто-то из детей перекладывает из своей тарелки по одной макаронине всем другим детям. После какого наименьшего количества действий у всех в тарелках может оказаться разное количество макаронин?

\begin{flushright}
        \scriptsize{5 баллов}
\end{flushright}

\problem{2}
В каждой клетке доски $8 \times 8$ написали по одному натуральному числу. Оказалось, что при любом разрезании доски на доминошки суммы чисел во всех доминошках будут разные. Может ли оказаться, что наибольшее записанное на доске число не больше 32 ? (Доминошкой называется прямоугольник, состоящий из двух клеток.)

\begin{flushright}
        \scriptsize{5 баллов}
\end{flushright}

\problem{3}
Произвольный треугольник разрезали на равные треугольники прямыми, параллельными сторонам (как показано на рисунке). Докажите, что ортоцентры шести закрашенных треугольников лежат на одной окружности. (Ортоцентр треугольника - точка пересечения его высот.)

\begin{flushright}
        \scriptsize{6 баллов}
\end{flushright}

\problem{4}
Квадратная коробка конфет разбита на 49 равных квадратных ячеек. В каждой ячейке лежит шоколадная конфета - либо чёрная, либо белая. Саша может съесть две конфеты, если они одного цвета и лежат в соседних по стороне или по углу ячейках. Какое наибольшее количество конфет гарантированно может съесть Саша, как бы ни лежали конфеты в коробке?


\begin{flushright}
        \scriptsize{8 баллов}
\end{flushright}

\problem{5}
На трёх красных и трёх синих карточках написаны шесть положительных чисел, все они различны. Известно, что на карточках какого-то одного цвета написаны попарные суммы каких-то трёх чисел, а на карточках другого цвета - попарные произведения тех же трёх чисел. Всегда ли можно гарантированно определить эти три числа?

\begin{flushright}
        \scriptsize{8 баллов}
\end{flushright}

\newpage

\hspace{3mm}

\problem{6}
Дан правильный $2 n$-угольник $A_1 A_2 \ldots A_{2 n}$ с центром $O$, причём $n \geq 5$. Диагонали $A_2 A_{n-1}$ и $A_3 A_n$ пересекаются в точке $F$, а $A_1 A_3$ и $A_2 A_{2 n-2}$ - в точке $P$. Докажите, что $P F=P O$.


\begin{flushright}
        \scriptsize{9 баллов}
\end{flushright}

\problem{7}
a) Группа людей прошла опрос, состоящий из 20 вопросов, на каждый из которых возможно два ответа. После опроса оказалось, что для любых 10 вопросов и любой комбинации ответов на эти вопросы существует человек, давший именно эти ответы на эти вопросы. Обязательно ли найдутся два человека, у которых ответы ни на один вопрос не совпали?
б) Решите ту же задачу, если на каждый вопрос есть 12 вариантов ответа.

\begin{flushright}
        \scriptsize{5+6 баллов}
\end{flushright}